\section{Числовые характеристики случайных величин.}
Пускай у нас всюду в этой секции задано некоторое вероятностное пространство $(\Omega, \mathcal{F}, \P)$.
\subsection{Математическое ожидание как интеграл Лебега.}
\subsubsection{Основные свойства.}
\begin{definition}
    Пусть $\xi: \Omega \to \R$ -- случайная величина.
    Будем называть математическим ожиданием случайной величины (и обозначать $\Ex \xi$) интеграл Лебега по $\P$ на $\Omega$:
    \[
        \Ex \xi = \int_{\Omega} \xi d\P.
    \]
    (естественно, математического ожидания может не существовать).
\end{definition}
Назовём несколько основных свойств математического ожидания:
\begin{enumerate}
    \item $\Ex c = c, \Ex{\Ind_A} = \P(A)$.
    \item $\Ex \alpha \xi + \beta \eta = \alpha \Ex \xi + \beta \Ex \eta$.
    \item Если две случайные величины $\xi$ и $\eta$ совпадают почти всюду, то $\Ex \xi = \Ex \eta$.
    \item Если $\xi \leq \eta$ почти всюду, то и $\Ex \xi \leq \Ex \eta$.
    \item Если $\xi \geq 0$ почти всюду, но $\Ex \xi = 0$, то $\xi = 0$ почти всюду.
    \item Пусть $\eta = \phi(\xi)$, где $\xi$ -- случайная величина, а $\phi$ -- борелевская функция.
    Тогда $\Ex \eta = \int_{\Omega} \phi(\xi)d\P = \Ex \phi(\xi)$.
    \item $|\Ex \xi| \leq \Ex |\xi|$.
\end{enumerate}
\subsubsection{Неравенства Йенсена и Ляпунова.}
\begin{theorem}[Неравенство Минковского]
    Для случайных величин $\xi$ и $\eta$ и некоторого $r > 1$ верно следующее:
    \[
        \|f + g\|_{r} \leq \|f\|_{r} + \|g\|_{r}.
    \]
\end{theorem}
\begin{theorem}[Неравенство Гёльдера]
    Для случайных величин $f \in L_p$ и $g \in L_{p'}$ справедливо следующее неравенство:
    \[
        \|f \cdot g\|_{1} \leq \|f\|_{p} \cdot \|g\|_{p'}.
    \]
\end{theorem}
\begin{theorem}[Неравенство Ляпунова]
    Для случайной величины $\xi$ и $u > v > 0$ справедливо следующее:
    \[
        \|\xi\|_{u} \geq \|\xi\|_{v}.
    \]
\end{theorem}
\begin{proof}
    Применим неравенство Гёльдера к $|\xi|^v \cdot 1$ и $u/v$. Тогда мы получим:
    \[
        \|\xi^v \cdot 1\|_1 \leq \|\xi^v\|_{u/v} \cdot  \cdot \|1\|_{p'} = \|\xi^v\|_{u/v}.
    \]
    Теперь возведём обе части неравенства в степень $1/v$ и получим:
    \[
        \|\xi\|_v = \biggr(\int |\xi|^v d\P\biggr)^{1/v} = (\|\xi^v\|_1)^{1/v} \leq \|\xi^v\|_{u/v}^{1/v} = \biggr(\int |\xi|^{u}\biggr)^{1/u} = \|\xi\|_u.
    \]
    Что и требовалось доказать.
\end{proof}
\begin{theorem}[Неравенство Йенсена]
    Для выпуклой вниз борелевской функции $g$ и случайной величины $\xi$ справедливо следующее неравенство:
    \[
        \Ex g(\xi) \geq g\bigr(\Ex \xi \bigr).
    \]
    (Если, естественно существуют $\Ex|g(\xi)|$ и $\Ex|\xi|$)
\end{theorem}
\begin{proof}
    Из выпуклости функции $g$ следует существование функции $d(x)$ такой, что
    \[
        g(y) \geq g(x) + d(x)(y - x).
    \]
    Пусть теперь $y = \xi$ и $x = \Ex \xi$. Возьмём математическое ожидание от обеих частей неравенства:
    \[
        \Ex g(\xi) \geq \Ex g(\Ex \xi) + \Ex \biggr( d(\Ex \xi)(\xi - \Ex \xi) \biggr) = g(\Ex \xi) + d(\Ex \xi) (\Ex \xi - \Ex \xi) = g(\Ex \xi).
    \]
\end{proof}
\subsubsection{Математическое ожидание независимых и некоррелированых случайных величин.}
\begin{theorem}
    Пусть $\xi$ и $\eta$ -- независимые случайные величины.
    Тогда, в случае существования у них первого момента, справедливо следующее равенство:
    \[
        \Ex \xi \eta = \Ex \xi \Ex \eta.
    \]
\end{theorem}
\begin{proof}
    Для простых функций данное утверждение очевидно.
    И действительно, если $\xi = \sum\limits_{k = 1}^{n} x_k I_{A_k}$ и $\eta = \sum\limits_{r = 1}^{m} y_k I_{B_r}$, то
    \[
        \Ex \xi \eta = \sum\limits_{k, r = 1, 1}^{n, m} x_k y_r \P(A_k B_r) = \sum\limits_{k, r = 1, 1}^{n, m} x_k y_r \P(A_k) \P(B_r) = \sum\limits_{k = 1}^n x_k \P(A_k) \cdot \sum\limits_{r = 1}^m y_r \P(B_r) = \Ex \xi \cdot \Ex \eta.
    \]
    В общем случае, для каждой неотрицательной случайной величины существуют последовательности $\xi_n$ и $\eta_n$, которые поточечно сходятся к $\xi$ и $\eta$ (без ограничения общности можно считать сходимость монотонной).
    Тогда, $\Ex \xi_n \eta_n =  \Ex \xi_n \Ex \eta_n$ -- по предыдущему пункту, и, совершая предельный переход по теореме Леви, получаем утверждение для неотрицательных функций. \\
    В общем случае для интегрируемых функций можно воспользоваться теоремой Лебега о мажорируемой сходимости и получить тот же результат или просто разбить на части $\xi$ и $\eta$ по знакам и для каждой из частей применить утверждение для неотрицательных случайных величин.
\end{proof}
\subsubsection{Экспоненциальное неравенство Маркова, неравенство Чернова.}
\begin{theorem}[Неравенство Маркова]
    Пусть $X$ -- неотрицательная случайная величина и $\epsilon > 0$. Тогда
    \[
        \P(X \geq \epsilon) \leq \dfrac{\Ex X}{\epsilon}.
    \]
\end{theorem}
\begin{proof}
    \[
        \P(X \geq \epsilon) = \int_{\{w \mid X(w) \geq \epsilon\}} d\P = \dfrac{1}{\epsilon} \int_{\{w \mid X(w) \geq \epsilon\}} \epsilon d\P \leq \dfrac{1}{\epsilon} \int_{\{w \mid X(w) \geq \epsilon\}} Xd\P \leq \dfrac{\Ex X}{\epsilon}.
    \]
\end{proof}
\begin{theorem}[Обобщённое неравенство Маркова]
    Пусть $\phi: \R \to (0, +\infty)$ -- строго возрастающая функция и $\epsilon \in \R$. Пусть $X$ -- случайная величина. Тогда
    \[
        \P(X \geq \epsilon) \leq \dfrac{\Ex \phi(X)}{\phi(\epsilon)}.
    \]
\end{theorem}
\begin{proof}
    Заметим, что так как $\phi$ -- строго возрастающая функция, то 
    \[
        X \geq \epsilon \Longleftrightarrow \phi(X) \geq \phi(\epsilon).
    \]
    А значит и 
    \[
        \P(X \geq \epsilon) = \P(\phi(X) \geq \phi(\epsilon)).
    \]
    Применим неравенство Маркова к $\phi(X)$ и получим искомую оценку:
    \[
        \P(X \geq \epsilon) \leq \dfrac{\Ex \phi(X)}{\phi(\epsilon)}.
    \]
\end{proof}
\begin{corollary}[Экспоненциальное неравенство Маркова]
    Пусть $X$ -- случайная величина, $\epsilon \in \R$ и $t > 0$. Тогда
    \[
        \P(X \geq \epsilon) \leq \dfrac{\Ex \exp(tX)}{\exp(t\epsilon)}.
    \]
\end{corollary}
\begin{proof}
    Применим обобщённое неравенство Маркова при $\phi(x) = e^{tx}$.
\end{proof}
\begin{corollary}[Неравенство Чернова]
    Пусть $X$ -- случайная величина и $a \in \R$. Тогда справедливо следующее неравенство:
    \[
        \P(X \geq a) \leq \inf\limits_{t > 0} e^{-ta}\Ex e^{tX}.
    \]
\end{corollary}
\begin{proof}
    Поскольку экспоненциальное неравенство Маркова для $X$ и $a \in \R$ справедливо при произвольном $t > 0$, то взятие $\inf$ по $t > 0$ завершает доказательство теоремы.
\end{proof}
\subsubsection{Неравенство Чебышёва для i.i.d.-последовательности.}
\begin{theorem}[Неравенство Чебышёва]
    Пусть $X$ -- случайная величина такая, что $\Ex |X| < +\infty$ и $\epsilon > 0$. Тогда верно следующее неравенство:
    \[
        \P(|X - \Ex X| \geq \epsilon) \leq \dfrac{\D X}{\epsilon^2}.
    \]
\end{theorem}
\begin{proof}
    Применим обобщённое неравенство Маркова при $\phi(x) = x^2$ и заметим, что $\Ex(X - \Ex X)^2 = \D X$.
\end{proof}
\begin{corollary}[Неравенство Чебышёва для i.i.d.-последовательности]
    Пусть $\{X_n\}$ -- последовательность независимых одинаково распределённых случайных величин с конечным первым моментом. Тогда для $S_n = \dfrac{1}{n}\sum\limits_{j = 1}^n X_j - \Ex X_j$ справедлива следующая оценка:
    \[
        \P(|S_n| \geq \epsilon) \leq \dfrac{\D X}{n\epsilon^2}.
    \]
\end{corollary}
\begin{proof}
    Применим неравенство Чебышёва к $S_n$ и $\epsilon$:
    \[
        \P(|S_n| \geq \epsilon) \leq \dfrac{\D S_n}{\epsilon^2}.
    \]
    Поскольку члены последовательности $S_n$ -- независимы, то 
    \[
        \D S_n = \dfrac{1}{n^2}\D\biggr(\sum\limits_{j = 1}^n X_j\biggr) = \dfrac{n}{n^2}\D X = \dfrac{\D X}{n}.
    \]
    Что и завершает доказательство теоремы.
\end{proof}
\subsubsection{Функции моментов, дисперсия, ковариация.}
\begin{definition}
    Для случайных величин $\xi$ и $\eta$ с конечным вторым моментом можно определить скалярное произведение следующим образом:
    \[
        (\xi, \eta) = \Ex \xi \eta.
    \]
\end{definition}
\begin{reminder}
    Пространство $L_2(\P)$ является гильбертовым пространство со скалярным произведением, определённым выше.
\end{reminder}
\begin{definition}
    Пусть $\xi$ -- случайная величина.
    Будем обозначать $\overset{\circ}{\xi}$ центрированную $\xi$, то есть $\overset{\circ}{\xi} = \xi - \Ex \xi$.
\end{definition}
\begin{definition}
    Будем называть $k$-м моментов математическое ожидание $\xi^k$, то есть $\Ex \xi^k$.
\end{definition}
\begin{definition}
    Будем называть $k$-м центральным моментом величину $\Ex \overset{\circ}{\xi} ^k$.
\end{definition}
\begin{definition}
    Будем называть $k$-м абсолютным моментом величину $\Ex |\xi|^k$
\end{definition}
\begin{definition}
    Второй центральный момент будем называть дисперсией и для случайной величины $\xi$ обозначать как
    \[
        \D \xi = \Ex (\xi - \Ex \xi)^2.
    \]
\end{definition}
\begin{property}
    Дисперсию можно переписать в немного другой форме:
    \[
        \D \xi = \Ex(\xi - \Ex \xi)^2 = \Ex \xi^2 - 2 \Ex (\xi \Ex \xi) + (\Ex \xi)^2 = \Ex \xi^2 - (\Ex \xi)^2.
    \]
\end{property}
\begin{property}
    Если $\eta = a\xi + b$, то 
    \[
        \D \eta = \Ex(\eta - \Ex \eta)^2 = \Ex(a\xi + b - a\Ex \xi - b)^2 = a^2 \Ex(\xi - \Ex \xi)^2 = a^2 \D \xi.
    \]
\end{property}
\begin{definition}
    Ковариацией случайных величин $\xi$ и $\eta$ будем называть 
    \[
        \cov(\xi, \eta) = \Ex(\overset{\circ}{\xi}\overset{\circ}{\eta}).
    \]
\end{definition}
\begin{property}
    Для случайных величин $\xi$ и $\eta$ ковариация выражается как
    \[
        \cov(\xi, \eta) = \Ex(\xi\eta) - \Ex\xi\Ex\eta.
    \]  
\end{property}
\begin{property}
    \[
        \cov(\xi, \xi) = \D\xi.
    \]
\end{property}
\begin{property}
    В силу неравенства Коши-Буняковского
    \[
        \cov^2(\xi, \eta) \leq \D\xi\D\eta.
    \]
\end{property}
\begin{property}
    Дисперсия суммы независимых случайных величин $\xi$ и $\eta$ равна сумме дисперсий:
    \[
        \D(\xi + \eta) = \D \xi + \D \eta.
    \]
\end{property}
\begin{proof}
    Распишем дисперсию по определению:
    \[
        \D(\xi + \eta) = \Ex\biggr(\overset{\circ}{\xi} + \overset{\circ}{\eta}\biggr)^2 = \Ex \overset{\circ}{\xi}^2  + 2\Ex \overset{\circ}{\xi}\overset{\circ}{\eta} + \Ex\overset{\circ}{\eta}^2 = \D \xi + 2 \cov(\xi, \eta) + \D \eta = \D \xi + \D \eta,
    \]
    где $\cov(\xi, \eta) = \Ex \overset{\circ}{\xi}\overset{\circ}{\eta} = \Ex \overset{\circ}{\xi}\Ex \overset{\circ}{\eta} = 0$ в силу независимости $\xi$ и $\eta$.
\end{proof}
\begin{property}
    Дисперсия суммы некоррелированных случайных величин также равно сумме дисперсий.
\end{property}
\subsubsection{Корреляция и её свойства.}
\begin{definition}
    Определим коэффициент корреляции для случайных величин $\xi$ и $\eta$ как 
    \[
        \rho(\xi, \eta) = \dfrac{\cov(\xi, \eta)}{\sqrt{\D \xi \D\eta}}.
    \]
\end{definition}
\begin{property}
    \[
    |\rho(\xi, \eta)| \leq 1.
    \]
\end{property}
\begin{property}
    Если $\xi$ и $\eta$ -- независимые случайные величины, то
    \[
        \rho(\xi, \eta) = 0.
    \]
\end{property}
\begin{property}
    Если $\eta = a \xi + b$, то
    \[
        \rho(\xi, \eta) = \dfrac{\cov(\xi, a \xi + b)}{\sqrt{\D\xi \D(a\xi + b)}} = \dfrac{a\cov(\xi, \xi)}{\sqrt{a^2\D \xi \D \xi}} = \sign a.
    \]
\end{property}